\section{Lattice support, projective coordinates, and cycle invariants}
\label{sec:lattice-coordinate-geometry}
Let $L$ be a nontrivial bounded distributive lattice and $K$ a field. Write
\[
 G_L(K)=\{(0,\lambda):\lambda\in L\}\cup\{(a,1):a\in K\}
 \subseteq K\times L,
\]
where addition uses join in the second coordinate and multiplication uses
meet. The repeated element $(0,1)$ is included only once. These operations
are closed: a nonzero amplitude can occur only at top support. The new
zero is $(0,0)$, and the unit is $(1,1)$. Thus this is a semiring, with
amplitude and support homomorphisms to $K$ and $L$. When $L=\B$ it is
$K_0$. We call a vector unimodular when its entries generate the unit ideal,
and take projective coordinates to mean its orbit under common scalar
units. This specifies an orbit functor without asserting a projective
scheme realization.

\begin{proposition}[Support profiles]\label{prop:lattice-profiles}
A vector of length $N$ over $G_L(K)$ is unimodular exactly when at least
one amplitude is nonzero. Its scalar orbit is specified by a nonempty
set $A\subseteq\{1,\ldots,N\}$ of top-support positions, a point of
$\mathbf P^{|A|-1}(K)$, and an arbitrary proper lattice element at every
position outside $A$. If $|L|=\ell<\infty$ and $K=\F_q$, the orbit count is
\[
 C_{N,L}(q)=\frac{(q+\ell-1)^N-\ell^N}{q-1}
 =\sum_{j=0}^{N-1}a_jq^j,
 \quad
 a_j=\sum_{r=j+1}^{N}\binom Nr(\ell-1)^{N-r}.
\]
Its formal point-counting zeta series is
\[
 \exp\!\left(\sum_{m\ge1}C_{N,L}(q^m)\frac{T^m}{m}\right)
 =\prod_{j=0}^{N-1}(1-q^jT)^{-a_j}.
\]
\end{proposition}
\begin{proof}
A coordinate of nonzero amplitude is a unit, so suffices for unimodularity.
If all amplitudes vanish, every linear combination has amplitude zero
and cannot be the unit. Fixing the top-support set leaves a nonzero vector
of $K^A$, on which $K^\times$ acts by the usual scalar action, and leaves
the proper lattice decorations fixed. The action on that nonzero vector
is free. A set $A$ of size $r$ therefore contributes
$(\ell-1)^{N-r}(q^r-1)/(q-1)$. Sum over $A$, and expand the geometric
sum to obtain the coefficients $a_j$. The last identity follows from
$\sum_{m\ge1}z^m/m=-\log(1-z)$ as a formal power series.
\end{proof}

\begin{proposition}[Linear symmetry and cycles]\label{prop:lattice-cycles}
For every nonzero commutative ring $R$,
\[
 \operatorname{GL}_n(G_L(R))=(R^\times)^n\rtimes\mathfrak S_n.
\]
For a matrix $A$ over $G_L(R)$, the signed Leibniz determinant has ordinary
determinant amplitude and support equal to the lattice permanent of its
support matrix. The $k$ th characteristic coefficient has support the join
of the support products of all principal cycle covers on $k$ vertices.
For a word in several matrices, its trace has support the join of the
corresponding colored closed-walk products. All these expressions are
invariant under monomial simultaneous conjugation.

Over a field, conjugacy classes of invertible monomial matrices under monomial conjugation are classified by the multiset
of permutation-cycle lengths and cycle products. A cycle of length $r$
and product $w$ contributes $t^r-w$ to the characteristic polynomial and
contributes $r w^{m/r}$ to the trace of the $m$ th power when $r\mid m$,
and zero otherwise. Block replication $A\mapsto A\otimes I_d$ raises
the split determinant and characteristic polynomial to the $d$ th power
and multiplies trace powers by the old scalar $d$.
\end{proposition}
\begin{proof}
Let $AB=BA=I$. In each diagonal amplitude sum some term is nonzero,
since the sum is $1$. Its two factors, say $A_{ik},B_{ki}$, have top support.
The off-diagonal support equations in $AB$ force $B_{kj}=0$ for $j\ne i$;
those in $BA$ force $A_{i l}=0$ for $l\ne k$. Here $0$ means the new zero.
The remaining diagonal equation says $A_{ik}B_{ki}=1$. Distinct rows have
distinct such columns, proving the monomial description. The converse is
immediate. If $R$ is the zero ring, this argument does not apply: a product
Boolean lattice can have different permutations on its different factors.

Apply amplitude and support to the Leibniz expansion. Signs are old units,
so have top support. This gives the determinant and principal-minor
formulas. Expanding a trace word lists exactly its colored closed walks.
A diagonal change of frame supplies a ratio of vertex units to each edge;
these ratios cancel on every cycle. A permutation relabels vertices,
proving invariance.

For conjugacy, a permutation change of basis identifies cycles of equal
length. Along a cycle choose one vertex scalar and successively determine
the others to match edge weights. The last edge agrees exactly when the
two cycle products agree. Ordering a basis by cycles gives the stated
characteristic and trace formulas directly. Finally a permutation of the
tensor-product basis turns $A\otimes I_d$ into $d$ diagonal copies of $A$,
with the new zero off the blocks. Matchings cannot cross those blocks,
so determinants and characteristic polynomials multiply, while traces add.
\end{proof}

A lattice join need not be an attained matching level. In $L=\B^2$, give
the diagonal entries of a $2$ by $2$ zero-amplitude matrix support $(1,0)$
and the off-diagonal entries support $(0,1)$. The two matching levels join
to $1$, although neither matching has level $1$. For Boolean support the
join does express existence of a matching; for a general lattice different
Boolean branches may have different witnesses.

\section{Measured adelic supports and logarithmic descent}
\label{sec:measured-defects}
A measured adelic family consists of a field $K$, a measure space
$(\Omega,\mathcal A,\nu)$, and absolute values $|\cdot|_\omega$ such that
$\ell_a(\omega)=\log|a|_\omega$ belongs to $L^1$ for every $a\in K^\times$.
It is proper when every $\ell_a$ has integral zero. This is the
Chen--Moriwaki framework \citep[Section 3.1]{ChenMoriwakiAdelic}. Let
$\mathcal B=\mathcal A/\text{null}$ be its Boolean measure algebra.
Define
\[
 \Delta([A])(a)=\int_A\ell_a\,d\nu,\qquad
 \mathcal D(f)(a)=\int_\Omega f\ell_a\,d\nu
 \quad(f\in L^\infty).
\]
Call $[A]$ balanced when $\Delta([A])=0$.

\begin{theorem}[Defect transform and localization]\label{thm:measured-defect}
The map $\Delta$ is a valuation from $\mathcal B$ to the additive group
$\operatorname{Hom}(K^\times,\mathbb R)$:
\[
 \Delta(A\cup B)+\Delta(A\cap B)=\Delta(A)+\Delta(B).
\]
The map $\mathcal D$ is real-linear, and $\mathcal D(1_A)=\Delta(A)$.
For a proper family, constants lie in its kernel and
$\Delta(A^c)=-\Delta(A)$. The bilinear pairing
\[
 \left(L^\infty/\ker\mathcal D\right)
 \times\operatorname{span}_{\mathbb R}\{\ell_a:a\in K^\times\}
 \longrightarrow\mathbb R,
 \qquad([f],g)\longmapsto\int fg\,d\nu
\]
is nondegenerate in both variables.

There is a canonical localization isomorphism
\[
 G_{\mathcal B}(K)[(0,[A])^{-1}]
 \cong\downarrow[A]\cong\mathcal B_{\Omega|A}.
\]
The restricted adelic family is proper exactly when $[A]$ is balanced.
A measure-preserving refinement preserving the absolute values commutes
with these defect maps and localization identifications.
\end{theorem}
\begin{proof}
Null changes do not affect integrals, and $\ell_{ab}=\ell_a+\ell_b$ proves
the character assertion. The indicator identity for union and intersection
proves the valuation formula. Bounded functions multiply integrable
functions into $L^1$, so $\mathcal D$ is well defined and linear. Properness
and $1_{A^c}=1-1_A$ give the remaining formulas.

The pairing is well defined by the definition of the kernel. A class
pairing trivially with every $\ell_a$ is zero. Conversely a nonzero
$g\in L^1$ pairs positively with the bounded function $\operatorname{sgn}g$,
since $\int|g|>0$. This proves nondegeneracy without claiming identification
with a full continuous or algebraic dual.

For any semiring and multiplicative idempotent $z$, localization at $z$
is $zS$, with unit $z$: every map inverting $z$ sends it to $1$ and factors
uniquely through $s\mapsto zs$. Here multiplication by $(0,[A])$ erases
amplitude and intersects support with $[A]$, giving $\downarrow[A]$.
Measurable subsets of $A$, modulo their restricted null sets, are exactly
that relative Boolean algebra. Its product formula is precisely
$\Delta(A)=0$. For a refinement $\pi$, change of variables gives
$\mathcal D'(f\circ\pi)=\mathcal D(f)$ and inverse images commute with
intersection, proving both naturality statements.
\end{proof}

The quotient in this theorem is a vector space. If the family is proper
and $\mathcal D\ne0$, its kernel contains $1$ but is proper, so it is not an
ideal of $L^\infty$. No quotient algebra or induced Boolean idempotent
structure is asserted for that vector-space quotient.

\begin{proposition}[Heights, neutral places, and metric states]
\label{prop:height-neutrality}
For $f\in L^\infty$ and $\mathbf x\in K^{n+1}\setminus\{0\}$, put
\[
 h_f(\mathbf x)=\int f(\omega)\log\max_i|x_i|_\omega\,d\nu.
\]
This is finite, and
$h_f(a\mathbf x)-h_f(\mathbf x)=\mathcal D(f)(a)$.
Consequently it descends to $\mathbf P^n(K)$ exactly when
$f\in\ker\mathcal D$. The same assertion, with the opposite sign,
holds for $-\int f\log\|s\|_\omega\,d\nu$ whenever the norm logarithms
of nonzero global vectors are integrable.

Adjoining any measured family $T$ of trivial absolute values leaves these
heights and canonical trivial-norm degrees unchanged. It replaces
$\mathcal B$ by $\mathcal B\times\mathcal B_T$ and the defect kernel by
$\ker\mathcal D\oplus_\infty L^\infty(T)$. In particular, for the standard
global place system, balanced supports after this extension are exactly
an empty or full original place set together with an arbitrary added
support.

Up to positive rescaling, the local states consisting of absence, a zero
seminorm on a line, and a nondegenerate seminorm form the chain
$0<1<2$, whose tensor product is meet. Measurable such states modulo null
sets form the bounded distributive lattice
\[
 \{(P,N)\in\mathcal B^2:N\le P\},
\]
where $P$ records presence and $N$ nondegeneracy. The function
$(P,N)\mapsto\Delta(N)$ is a valuation on this lattice.
\end{proposition}
\begin{proof}
Discard zero coordinates. The logarithm of the maximum is the maximum of
finitely many $L^1$ functions, and its absolute value is bounded by the sum
of their absolute values. Scalar multiplication adds $\ell_a$ to that
maximum. The norm formula follows in exactly the same way from
$\|as\|_\omega=|a|_\omega\|s\|_\omega$. Invariance under every scalar
is equivalent to vanishing of the displayed character, including when
$n=0$.

On the added family every principal logarithm is zero and the maximum
of the absolute values of a nonzero tuple is $1$. The canonical trivial
norm is $1$ on nonzero vectors. Thus all additional integrands vanish.
Measurable sets and essentially bounded functions on a disjoint union
split into their two components, yielding the product and kernel formulas.
The global-field assertion follows from Theorem~\ref{thm:global-weights}.

A seminorm on a real line is either zero or a positive multiple of any
fixed norm. Tensoring with absence remains absent, tensoring a present
zero seminorm gives a zero seminorm, and tensoring two nondegenerate
seminorms remains nondegenerate. These are the three meet rules. The
threshold maps $j\mapsto1_{j\ge1}$ and $j\mapsto1_{j\ge2}$ identify the
chain with the incidence pairs in $\B^2$. Apply this pointwise; the inverse
assigns $2$ on $N$, $1$ on $P\setminus N$, and $0$ outside $P$. Projection to $N$
preserves joins and meets, so the valuation identity follows. This
incidence lattice is generally not Boolean.
\end{proof}

Booleanization of an absolute value itself only tests whether a scalar
vanishes: $|a|_\omega>0$ for every $a\ne0$. It is therefore independent of
the place. A trivial absolute value is nondegenerate and differs from an
identically zero seminorm.

The measure algebra above is always Boolean sigma-complete, but need not
be complete. To see this, equip an uncountable set $X$ with its
countable/cocountable sigma-algebra and counting measure, and use the
trivial absolute value at every point. All principal logarithms vanish,
so the family is proper. Its null ideal is $\{\varnothing\}$. If $Y$ and
$X\setminus Y$ are both uncountable, the singleton subsets of $Y$ have no
least measurable upper bound: any such upper bound is cocountable and
can be made strictly smaller by deleting a point outside $Y$. Thus the
unqualified completeness claim fails. The displayed constructions use
only finite lattice operations and require no completeness assumption.

\section{Monomial bundles and support-refined invariant gluing}
\label{sec:monomial-invariant-gluing}
\begin{proposition}\label{prop:cover-local-system}
In the usual local-system category of a connected space, rank $n$ local
systems with monomial transition maps are equivalent to degree $n$ finite
covering local systems equipped with a rank-one local system on the cover.
For coefficients $\mathbb C_0$ these are exactly the local systems defined
by invertible $n$ by $n$ split matrices. Their ordinary loop determinants
factor over the cycles of the covering permutation.
\end{proposition}
\begin{proof}
The permutation parts transport the $n$ basis labels and define the cover.
On each label the diagonal multiplier transports a one-dimensional fiber,
defining the line local system. Conversely the direct sum over the sheets
has monomial transition maps. The constructions are inverse on objects
and on their compatible isomorphisms. The split assertion is the general
linear-group calculation, and the loop determinant is the cycle formula
in Proposition~\ref{prop:lattice-cycles}.
\end{proof}

For a sheaf of modules or line bundle $V$, write $\operatorname{Spl}(V)$ for
the sheaf coproduct of $V$ and an absent point. A section is locally either
an old section or absent; addition uses the old sum when both are present,
and absence is its identity. Products take values in the corresponding
tensor product and are absent if either factor is absent.

\begin{proposition}[Matrix Higgs data]\label{prop:matrix-higgs-support}
Let $X$ be a smooth complex curve. Consider a vector bundle presented by
monomial transition matrices, together with local $n$ by $n$ matrices of
sections of $\operatorname{Spl}(K_X)$ that transform by conjugation.
The split characteristic coefficients define sections
\[
 c_k\in\Gamma(X,\operatorname{Spl}(K_X^{\otimes k})),
 \qquad1\le k\le n.
\]
Their amplitudes are the ordinary characteristic coefficients. At a
geometric point $c_k$ is absent exactly when no principal $k$-vertex cycle
cover is present; it is the old zero exactly when some such cover is
present and the amplitude coefficient vanishes.
\end{proposition}
\begin{proof}
Expand $c_k$ as the signed sum of its principal cycle-cover products. Each
product has tensor degree $k$. A diagonal change of bundle frame cancels
around each cycle; a permutation relabels the covers. Thus the local
expressions agree on overlaps in $K_X^{\otimes k}$ and glue in its split
sheaf. Applying amplitude and Boolean support to the expansion gives the
two assertions at a geometric point.
\end{proof}

These are arbitrary matrix Higgs fields on a monomially reduced bundle.
They are not restricted to the Lie algebra of the monomial normalizer,
which would impose diagonal fields. A further restriction of a geometric
Langlands functor requires a specified spectral stack, functor and all of
its hypotheses; the invariant-gluing proposition supplies none of those
additional structures.

\begin{proposition}[A split torsor calculation]\label{prop:split-torsor}
Let $M$ be a nonempty simply transitive $W$-set, and adjoin two distinct
$W$-fixed points. Its action groupoid is equivalent to
$BW\sqcup BW\sqcup *$. Finite products give the corresponding product
groupoids and a Boolean support map recording absence at each factor.
\end{proposition}
\begin{proof}
The two fixed orbits each have stabilizer $W$ and give $BW$. A simply
transitive orbit has trivial stabilizer, so its action groupoid is
contractible. The decomposition is invariant under the action, and
products of these decompositions prove the final assertion. Applying this
calculation to a proposed arithmetic torsor still requires that torsor to
have been defined and its simple transitivity proved.
\end{proof}
